% Nouveau EGA V (cote 25): archive cote 25, site page 163 (PDF page 164).
% typescript. Draft compared: none.
% AI-transcribed by Claude Opus 5.5 (claude-opus-5-5, adaptive thinking, effort high), Message Batches API, batch ega25b (msgbatch_01A7n5YuXubUPNK1zsepi8RG), 2026-11-13; not yet checked by eye.
\scanpage{25/163}
\mspage{\hand{6}}

un sous-préschéma X" de X, régulier, \xout{tel que} contenant X, et tel que X
puisse se décrire dans X"\textsubscript{\xout{\unsure{w}}} par une seule équation $\hand{\phi} = 0$ , admettant x comme
zéro \inserted[typed]{sing} élémentaire de multiplicité n.

\inserted[typed]{On conclut en particulier que si A est local noeth élémentaire de}
\xout{15.7. On peut donner des variantes géométriques de 15.6\unsure{.,|,} \unsure{lorsqu'on}}
\inserted[typed]{mult. n, alors A est Coh. Mac. \xout{Il est clair} Nous montrerons au par.20}
\xout{est sur un corps k (point de X, pas nécessairement rationnel sur k)}
(par un calcul facile d'éclatement) que le point fermé de Spec(A) est
le seul point singulier de Spec(A); il en résulte que \xout{si dim} A est normal
sss dim A $\hand{\geqslant}$ 2, réduit sss dimA $\hand{\geqslant}$1. Si dim A = 0, alors  A est élémentaire
de multiplicité n sss  $m/m^2$ est de \over{dim}{rang}1 et  n est le plus petit entier
tel que $\underline{m}^n = 0$; on notera que ces anneaux sont aussi lesquotients d'anneaux
de valu\xout{\unsure{at}}ion discrète par la\inserted[typed]{\xout{\unsure{ent}}} puissance n.ème de \deleted{l'}\inserted[typed]{\unsure{l|\hand{l}}eur} idéal maximal.

15.7.   Il y a lieu de dégager les variantes ``géométriques'' de 15.6. ,
on trouvera en particulier que si X est un préschéma loc de type fini sur
le corps k, $x \hand{\in} X$, et si \over{n}{d} $= \dim_x X$ , alors  x est une singularité \xout{ordinaire}
\inserted[typed]{élémentaire}
\xout{\unsure{re}} de multiplicité  n de X, sss il existe un voisinage ouvert U de x
tel que U se plonge comme sous-préschéma d'un préschéma lisse X' sur k,
\xout{de d} conne xe et de dimension \over{n}{d}+1 , défini par une équation $\hand{\phi}$=0 , x étant
un zéro \inserted[typed]{géom.} singulier ordinai\unsure{r|m}e de multiplicité n de $\hand{\phi}$ . On aura dit d'ailleurs
\margin{$\surd$} dans 15.5. que cela signifie  que  la \inserted[typed]{valeur en x de} partie principale $d^n\hand{\phi}$ de $\hand{\phi}$ , \xout{restrei}
\xout{nte à X=V($\hand{\phi}$)} (qui à priori est un\xout{e section} \inserted[typed]{élément} de $\underline{P}^n_{X'/k}\hand{\otimes} \unsure{\underline{k}(x)}$ ,et plus
précisément de son Idéal d'augmentation $\underline{I}$ )  est en fait un\xout{e section} \inserted[typed]{élément de}
\xout{de} $\underline{I}^n = \underline{Sym}^n_{\underline{O}_U}(\hand{\underline{\Omega}}^1_{X'/k}\hand{\otimes} \unsure{k(\hand{x})}\textsubscript{\xout{\unsure{w}}})$), et  à ce titre est une \xout{section} \inserted[typed]{forme} lisse \xout{en x}
\margin{|} \hand{[}NB Oubli dans 15.1. : l'ensemble des points de lissité d'une forme est
\margin{|} ouvert\hand{]}. On remarquera de même qu'un tel point est isolé dans l'ensemble
des points non lisses, il est géom normal \xout{\illegible{1}} (resp. géom réduit) sss
\over{n}{d} $\hand{\geqslant} 2$  (resp. \over{n}{d} $\hand{\geqslant} 1$). Si \over{n}{d} = 0 et k(x) = k (\inserted[typed]{i.e. x  isolé et} x rationnel sur k) alors \xout{x}
point singulier élémentaire  de mult\hand{\unsure{,}} \unsure{n} signifie que $\underline{O}_{X,x}$  est  \deleted{k-isomor}

\mspage{\hand{163}}
